\section{Related work}
\label{sec:rw}

\myparagraph{Learning 3D {semantic} scene representations with NeRF}
NeRF~\citep{mildenhall2021nerf} uses a multilayer perceptron to predict the volume density and radiance for any given 3D position and viewing direction.
Such a representation can naturally be extended to semantic features.
The early works N3F~\citep{tschernezki2022n3f} and DFF~\citep{kobayashi2022decomposing} distill DINO 2D (\emph{i.e.}, image-level) features~\citep{caron2021dino} in scene-specific NeRF representations. \citet{kobayashi2022decomposing} also distill LSeg~\citep{li2022lseg}, a  language-driven model for semantic segmentation. Building on this, LERF~\citep{kerr2023lerf} and 3D-OVS~\citep{liu2023weakly} learn 3D CLIP~\citep{radford2021clip} and DINO~\citep{caron2021dino} features jointly for open-vocabulary segmentation. These works were extended to other pretrained models such as latent diffusion models~\citep{ye2023featurenerf} or  SAM~\citep{kirillov2023sam} for segmentation~\citep{cen2023segmentNeRF, ying2023omniseg3d}.

\myparagraph{Learning 3D {semantic} scene representations with GS}
Subsequent work has relied on the Gaussian Splatting method~\citep{kerbl2023gaussiansplatting}, enabling rendering speeds orders of magnitude faster than NeRF-based models. 
Several tasks have been addressed such as semantic segmentation using SAM~\citep{cen2023segmentGS, ye2024gaussiangrouping, kim2024garfield, wu2024opengaussian}, language-driven retrieval or editing using CLIP combined with DINO~\citep{zuo2024fmgs} or SAM~\citep{ye2023featurenerf, wu2024opengaussian}, scene editing using diffusion models~\citep{chen2024gaussianeditor, wang2024gaussianeditor}, and 3D-aware finetuning~\citep{yue2024fit3d}.
These works learn 3D semantic representations by minimizing a reprojection loss.
As a single scene can be represented by over a million Gaussians, such optimization-based techniques have strong memory and computational limitations. 
To handle these, FMGS~\citep{zuo2024fmgs} employs a multi-resolution hash embedding (MHE) of the scene for uplifting DINO and CLIP representations, Feature 3DGS~\citep{zhou2024feature3dgs} learns a $1\times 1$ convolutional upsampler of Gaussians' features distilled from LSeg and SAM's encoder, and LangSplat~\citep{qin2023langsplat} learns an autoencoder to reduce CLIP feature dimension from 512 to 3. In contrast, our approach requires no learning, which significantly speeds up the uplifting process and reduces the memory requirements. 

\myparagraph{Direct uplifting of 2D features into 3D} Direct uplifting from 2D to 3D has been explored in prior works.
GaussianEditor~\citep{chen2024gaussianeditor} uplifts 2D SAM masks into 3D to selectively optimize Gaussians for editing tasks, see details in the supplementary (Sec.~\ref{sec:appendx-gaussianeditor}). Dr. Splat~\citep{jun2025drsplat} also lifts 2D SAM masks, to create an assignment from object CLIP features to Gaussians. Both approaches are specific to segmentation and are ill-suited for uplifting generic representations. Semantic Gaussians \citep{guo2024semantic} pairs 2D pixels with 3D Gaussians along each pixel’s ray based on depth information but relies on a learned 3D convolutional network. OpenGaussian \citep{wu2024opengaussian} also pairs pixels with  Gaussians, but relies on pseudo-features learned with a contrastive loss on SAM masks. In contrast, our approach is simple to implement, applicable to any 2D representation and entirely parameter-free.

\myparagraph{Leveraging 3D information to better segment in 2D}
Most prior works focusing on semantic segmentation leverage 2D models specialized for this task. The early work of \cite{yen2022nerfsupervision} uplifts learned 2D image inpainters by optimizing view consistency over depth and appearance. 
Subsequent works have mostly relied on uplifting either features from SAM's encoder~\citep{zhou2024feature3dgs}, binary SAM masks~\citep{cen2023segmentNeRF, cen2023segmentGS}, or SAM masks automatically generated for all objects in the image~\citep{ye2024gaussiangrouping, ying2023omniseg3d, kim2024garfield}.
The latter approach is computationally expensive, as it requires querying SAM on a grid of points over the image. It also requires matching inconsistent mask predictions across views, with \eg a temporal propagation model~\citep{ye2024gaussiangrouping} or a hierarchical learning approach~\citep{kim2024garfield}, which introduces additional computational overhead. In this work, we focus on single instance segmentation and show that our uplifted features are on par with the state of the art~\citep{cen2023segmentNeRF,cen2023segmentGS,ying2023omniseg3d}.
Standing out from prior work uplifting DINO features~\citep{tschernezki2022n3f,kobayashi2022decomposing,isrfgoel2023interactive,kerr2023lerf,liu2023weakly, ye2023featurenerf, zuo2024fmgs}, we show that DINOv2 features can be used on their own for semantic segmentation and rival SAM-based models through a simple graph diffusion process that leverages 3D geometry.

\myparagraph{3D CLIP features for open-vocabulary localization} 
For learning 3D CLIP features, prior works also leverage vision models such as DINO or SAM. DINO is used to regularize and refine CLIP features~\citep{kerr2023lerf, liu2023weakly, zuo2024fmgs, shi2024legaussians}, while SAM is employed for generating instance-level CLIP representations~\citep{qin2023langsplat, wu2024opengaussian, jun2025drsplat}. These approaches suffer from high computational costs, resorting to dimensionality reduction or efficient multi-resolution embedding representations, and usually run for a total of one to two hours for feature map generation and 3D feature optimization. 
In contrast, our approach bypasses the high computational cost of gradient-based optimization and, combined with graph diffusion, is an order of magnitude faster than these prior works.